The rise of streaming platforms has increased saturation in the music industry, driving artists to seek new ways to differentiate themselves and connect with their audiences. In this context, virtual reality (VR) emerges as an alternative for creating immersive, interactive, and personalized musical experiences. However, most studies on musical experiences in VR focus on live concerts, active navigation, multiple perspectives, or social interaction between the audience and the performer. This leaves an opportunity to explore music video experiences, where visual, auditory, and interactive elements are organized through temporal sequencing designed to accompany the narrative and identity of a specific song.

This project proposes the development of an immersive sequencing system based on Event-Driven Architectures (EDA), which makes it possible to organize and coordinate visual, interactive, and auditory elements within a VR music video. Through timed events, changes in environment, visual stimuli, interactions, and sound elements are synchronized with the structure of the song, so that the experience remains coherent with its musical narrative.

The experience is evaluated using the I-Group Presence Questionnaire (IPQ), an instrument that measures the perception of presence through dimensions such as spatial presence, involvement, and realism. Through these dimensions, the project seeks to analyze whether event-based sequencing contributes to strengthening the user's perception of presence and offering a personalized, immersive VR experience aligned with the project's artistic intention.
